
\subsection{算子 $\mathfrak q$ 的构造 II：$\mathfrak q$ 是 $L_\infty$ 同态}
\label{subsec:7-4-3}

本小节把定理~\ref{thm:3-8-32}
的结论表述为滤过 $L_\infty$ 同态的语言。\liftIndex{I287}%
滤过 $L_\infty$ 代数\liftIndex{I286}%
及 $L_\infty$ 同态的定义见第 A3 节。

设 $(C,\mathfrak m)$ 为滤过 $A_\infty$ 代数，
$(C^{\mathrm{amb}},\delta)$ 为
$\Lambda_{0,\mathrm{nov}}$ 上的上链复形。
在定理~\ref{thm:3-8-32} 的证明中，
将取
$C^{\mathrm{amb}}
=C(M;\Lambda_{0,\mathrm{nov}})$、
$C=C(L;\Lambda_{0,\mathrm{nov}})$。
在 $C^{\mathrm{amb}}[1]$ 上定义平凡的
滤过 $L_\infty$ 结构：
\begin{equation}
\begin{cases}
\mathfrak l_1(y)=(-1)^{\deg y}\delta(y),\\
\mathfrak l_k=0, & k\ne1.
\end{cases}
\tag{7.4.20}\label{eq:7-4-20}
\end{equation}
为何这里取符号 $(-1)^{\deg y}$
而不是 $(-1)^{\deg' y}$，
我们没有很好的解释；
唯一的理由是取前者时
命题 7.4.34 的证明能够成立。

接下来考虑 Hochschild 复形：
\begin{equation}
CH(C,C)=\prod_{k=0}^{\infty}
\operatorname{Hom}^{\mathrm{filt}}
\bigl(B_k(C[1]),C[1]\bigr).
\tag{7.4.21}\label{eq:7-4-21}
\end{equation}
这里 $\operatorname{Hom}^{\mathrm{filt}}$
由所有满足
\[
\varphi\bigl(F^\lambda(B_k(C[1]))\bigr)
\subset F^\lambda C[1]
\]
的 $\Lambda_{0,\mathrm{nov}}$-模同态
$\varphi$ 组成。
以下称这类同态为\emph{滤过同态}。

我们将在 $CH(C,C)$ 上如下定义
微分分次李代数（D.G.L.A.）结构。
分次由同态的次数决定，
而非式~\eqref{eq:7-4-21} 右端的 $k$。
定义微分和括号之前，
需要几个引理和一些记号。
以下讨论是第~\ref{subsec:4-4-5} 小节的滤过版本。

\begin{definition}{7.4.22}\label{def:7-4-22}
以
$\operatorname{Der}
(\widehat B(C[1]),\widehat B(C[1]))$
表示所有滤过同态
\[
\widehat\varphi:
\widehat B(C[1])\longrightarrow
\widehat B(C[1])
\]
所成的集合，要求它们满足

\par\noindent\textup{(7.4.22.1)}\label{cond:7-4-22-1}\quad
\[
\bigl((\widehat\varphi\widehat\otimes
\operatorname{id})
+(\operatorname{id}\widehat\otimes
\widehat\varphi)\bigr)\circ\Delta
=\Delta\circ\widehat\varphi;
\]

\par\noindent\textup{(7.4.22.2)}\label{cond:7-4-22-2}\quad
对任意
$\mathbf x\in\widehat B(C[1])$，
$\widehat\varphi(\mathbf x)$ 的
$B_0(C[1])$ 分量为零。
注意，若
$1\in\widehat B_0(C[1])$，
则 $\widehat\varphi(1)$ 仍可非零。
\end{definition}

回顾记号
\[
\Delta^{n-1}\mathbf x
=\sum_a
\mathbf x_a^{(n;1)}\otimes\cdots\otimes
\mathbf x_a^{(n;n)}.
\]
这里求和指标 $a$ 的取值集合
依赖于 $n$ 和 $\mathbf x$。
对每个 $\varphi=(\varphi_0,\varphi_1,\ldots)$，
定义
$\widehat\varphi:
\widehat B(C[1])\to\widehat B(C[1])$：
\begin{equation}
\widehat\varphi(\mathbf x)
=\sum_a
(-1)^{\deg\varphi\,\deg'\mathbf x_a^{(3;1)}}
\mathbf x_a^{(3;1)}
\otimes\varphi(\mathbf x_a^{(3;2)})
\otimes\mathbf x_a^{(3;3)}.
\tag{7.4.23}\label{eq:7-4-23}
\end{equation}
其中
$\varphi:\widehat B(C[1])\to C[1]$
是一个同态，
其在 $\widehat B_k(C[1])$ 上的限制
为 $\varphi_k$。

与引理~\ref{lem:4-4-43} 相同的方法
可证明下面的结论。

\begin{lemma}{7.4.24}\label{lem:7-4-24}
映射 $\varphi\mapsto\widehat\varphi$
定义同构
\begin{equation}
CH(C,C)\cong
\operatorname{Der}
(\widehat B(C[1]),\widehat B(C[1])).
\tag{7.4.25.1}\label{eq:7-4-25-1}
\end{equation}
若 $C$（严格）含幺，则定义
\begin{equation}
CH^{\mathrm{unit}}(C,C)\cong
\operatorname{Der}^{\mathrm{unit}}
(\widehat B(C[1]),\widehat B(C[1]))
\tag{7.4.25.2}\label{eq:7-4-25-2}
\end{equation}
为满足
\[
\varphi(\mathbf x_1\otimes\mathbf e\otimes
\mathbf x_2)=0
\]
的所有 $\varphi$ 所成的集合，
其中任取
$\mathbf x_1,\mathbf x_2\in
\widehat B(C[1])$。
\end{lemma}

以下经同构
\eqref{eq:7-4-25-1}，
把 $CH(C,C)$ 与
$\operatorname{Der}
(\widehat B(C[1]),\widehat B(C[1]))$
等同。另记
\[
\begin{aligned}
\operatorname{Der}B(C[1])
&:=\operatorname{Der}
(\widehat B(C[1]),\widehat B(C[1])),\\
\operatorname{Der}B^{\mathrm{unit}}(C[1])
&:=\operatorname{Der}^{\mathrm{unit}}
(\widehat B(C[1]),\widehat B(C[1])).
\end{aligned}
\]
